\section{\textcolor{titleblue}{Methods}}
\subsection{RNA design problem}
\subsubsection{Sequences, structures, and valid designs}
To approach the RNA design problem, we first define the valid design space
for a target structure. A sequence of length $n$ can be represented as
$x=x_1\cdots x_n$, where each nucleotide $x_i$ is chosen from $\A=\{\mathrm{A,C,G,U}\}$. Given a target structure
$y$, we can further restrict the base choices using two sets:
$\mathcal{P}(y)$ for paired positions and $\mathcal{U}(y)$ for unpaired positions~\citep{tang2026}. Each pair in $\mathcal{P}(y)$
must be one of the canonical or wobble pairs:
\begin{equation}
 \mathcal{C}=\{\mathrm{AU,UA,CG,GC,GU,UG}\}.
\end{equation}
The valid design space can then be represented as
\begin{equation}
 \begin{split}
 \X(y)=\{\,&x\in\A^n:\;x_ix_j\in\mathcal{C}\\
 &\text{for every }(i,j)\in\mathcal{P}(y)\}.
 \end{split}
 \label{eq:validspace}
\end{equation}
This restriction greatly reduces the design space for structures with many pairs by
removing incompatible choices, such as C pairing with A. The number of sequences
in the design space for a target $y$ with $p$ pairs and $u$
unpaired positions is then
\begin{equation}
 |\X(y)|=6^p4^u,\qquad |y|=2p+u.
\end{equation}

\subsubsection{Target folding probability}
Now that we have a valid design space for the target $y$,
we need a way to assess how well a given sequence favors folding
into the target. To do this, we calculate the target folding probability,
which is the probability that the sequence folds into the target structure,
given its ensemble of competing structures. Let $G(x,y)$ be the free energy of
a sequence $x$ folded into structure $y$, which is calculated using the Turner
2004 nearest-neighbor energy model~\citep{mathews2004}. The lower the free energy, the more favorable the
structure is for that sequence. The folding ensemble for sequence $x$ is
\begin{equation}
 \mathcal{Y}(x)=\{y': y'\text{ is an allowed structure for }x\}.
\end{equation}
This set contains the target structure and all competing folds that the sequence
can fold into. The total statistical weight of all structures for the sequence
is given by the partition function~\citep{mccaskill1990},
\begin{equation}
 Z(x)=\sum_{y'\in\mathcal{Y}(x)}e^{-G(x,y')/(RT)}.
 \label{eq:partition}
\end{equation}
Here $R$ is the gas constant and $T$ is the temperature;
the product $RT$ is usually expressed in kcal/mol. Based on this formula,
the lower the energy of a structure, the more it contributes to the ensemble weight.

Target folding probability can then be defined as
\begin{equation}
 P(y\mid x)=\frac{e^{-G(x,y)/(RT)}}{Z(x)}.
 \label{eq:probability}
\end{equation}
This equation divides the target's Boltzmann weight by the total weight of all
allowed structures, giving the probability that the sequence forms the target under the thermodynamic model.

The primary design objective of EliteDesign is
\begin{equation}
 x^\star\in\underset{x\in\X(y)}{\arg\max}\;P(y\mid x).
 \label{eq:design}
\end{equation}
During search we rank candidates by log probability,
\begin{equation}
 \ell(x)=\log P(y\mid x)
 =-\frac{G(x,y)}{RT}-\log Z(x),
 \label{eq:logprob}
\end{equation}
We optimize in log space to avoid underflow when probabilities are very small.

\subsubsection{Normalized ensemble defect}
Normalized ensemble defect (NED) is the expected fraction of nucleotides whose pairing state
differs from the target structure $y$, averaged over all folds $y'$ in the
ensemble~\citep{dirks2004,zadeh2011}. For positions $i$ and $j$, the pairing probability is
\begin{equation}
 p_{ij}(x)=\sum_{(i,j)\in\mathcal{P}(y')}P(y'\mid x).
\end{equation}
The probability that position $i$ is unpaired is
\begin{equation}
 u_i(x)=1-\sum_{j\ne i}p_{ij}(x).
\end{equation}
Let $\mathcal{P}$ and $\mathcal{U}$ be the sets containing the target's paired and unpaired positions respectively. Then
{\small
\begin{equation}
 \mathrm{NED}(x)=1-\frac{1}{n}\biggl[
 2\textstyle\sum_{(i,j)\in\mathcal{P}}p_{ij}(x)
 +\textstyle\sum_{i\in\mathcal{U}}u_i(x)\biggr].
 \label{eq:ned}
\end{equation}
}
Each correct pair is counted twice because it contains two nucleotides.

\subsection{Structure parsing and distribution initialization}
Following SamplingDesign~\citep{tang2026}, given a target structure in dot-bracket notation, for example
\[
\setlength{\arraycolsep}{3pt}
\begin{array}{cccccccccc}
 ( & . & ( & . & . & . & ) & . & . & )\\
 \scriptstyle 1 & \scriptstyle 2 & \scriptstyle 3 &
 \scriptstyle 4 & \scriptstyle 5 & \scriptstyle 6 &
 \scriptstyle 7 & \scriptstyle 8 & \scriptstyle 9 &
 \scriptstyle 10
\end{array}
\]
we parse the structure into four types of variables:
\begin{itemize}
\item \textbf{Unpaired variables} with four nucleotide choices: A, C, G, or U.
  For example, position 5 is an unpaired variable.
\item \textbf{Paired variables} sampled from {AU, UA, CG, GC, GU, or UG}.
  For example, positions $(1,10)$ form a paired variable.
\item \textbf{Mismatch variables} group the two unpaired positions next to the pair that closes
  a loop, and have $4^2=16$ choices. For example, positions $(4,6)$ form a mismatch variable.
\item \textbf{Trimismatch variables} group three unpaired positions in a loop, one on one side
  and two on the other side. These bases are chosen together,
  giving $4^3=64$ possible combinations. For example, positions $(2,8,9)$ form a trimismatch variable.
\end{itemize}

\begin{figure}[t]
\centering
\begin{tikzpicture}[
  x=1.22cm,y=1.22cm,
  base/.style={circle,fill=white,line width=1.4pt,
    minimum size=0.42cm,inner sep=1pt,font=\normalsize\bfseries,text=black},
  position/.style={font=\normalsize\bfseries,text=black}
]
\filldraw[fill=vartwocolor!12,draw=vartwocolor!65,line width=0.8pt]
  (-1,0.865)--(1.431,1.365)--(1.431,0.365)--cycle;
\draw[varfourcolor!28,line width=9pt,line cap=round]
  (-0.809,2.681)--(0.809,2.681);
\draw[black!45,line width=1pt]
  (-0.5,0)--(-1,0.865)--(-0.5,1.73)--(-0.809,2.681)--
  (0,3.269)--(0.809,2.681)--(0.5,1.73)--
  (1.431,1.365)--(1.431,0.365)--(0.5,0);
\draw[varonecolor,line width=2.2pt] (-0.5,0)--(0.5,0);
\draw[varthreecolor,line width=2.2pt] (-0.5,1.73)--(0.5,1.73);
\node[font=\normalsize\bfseries,text=black,fill=white,inner sep=1pt]
  at (0,-0.32) {$v_1$};
\node[font=\normalsize\bfseries,text=black,fill=white,inner sep=1pt]
  at (0,1.46) {$v_3$};
\node[font=\normalsize\bfseries,text=black] at (0.66,0.85) {$v_2$};
\node[font=\normalsize\bfseries,text=black] at (0,2.681) {$v_4$};
\node[base,draw=varonecolor] (n1) at (-0.5,0) {G};
\node[base,draw=vartwocolor] (n2) at (-1,0.865) {A};
\node[base,draw=varthreecolor] (n3) at (-0.5,1.73) {C};
\node[base,draw=varfourcolor] (n4) at (-0.809,2.681) {C};
\node[base,draw=varfivecolor] (n5) at (0,3.269) {G};
\node[base,draw=varfourcolor] (n6) at (0.809,2.681) {A};
\node[base,draw=varthreecolor] (n7) at (0.5,1.73) {G};
\node[base,draw=vartwocolor] (n8) at (1.431,1.365) {U};
\node[base,draw=vartwocolor] (n9) at (1.431,0.365) {C};
\node[base,draw=varonecolor] (n10) at (0.5,0) {C};
\node[position] at (-0.5,-0.37) {1};
\node[position] at (-1.37,0.865) {2};
\node[position] at (-0.88,1.7) {3};
\node[position] at (-1.18,2.72) {4};
\node[position] at (0,3.64) {5};
\node[position] at (1.18,2.72) {6};
\node[position] at (0.84,1.98) {7};
\node[position] at (1.8,1.365) {8};
\node[position] at (1.8,0.365) {9};
\node[position] at (0.5,-0.37) {10};
\node[font=\normalsize\bfseries] at (-1.02,-0.08) {$5'$};
\node[font=\normalsize\bfseries] at (1.02,-0.08) {$3'$};
\end{tikzpicture}
\smallskip

{\normalsize
\begin{tabular}{@{}c@{\,}r@{\quad}l@{\quad}l@{}}
\textcolor{varonecolor}{\rule{0.65em}{0.65em}} & $v_1$ & Pair & $(1,10)$\\
\textcolor{vartwocolor}{\rule{0.65em}{0.65em}} & $v_2$ & Trimismatch & $(2,8,9)$\\
\textcolor{varthreecolor}{\rule{0.65em}{0.65em}} & $v_3$ & Pair & $(3,7)$\\
\textcolor{varfourcolor}{\rule{0.65em}{0.65em}} & $v_4$ & Mismatch & $(4,6)$\\
\textcolor{varfivecolor}{\rule{0.65em}{0.65em}} & $v_5$ & Unpaired & $(5)$
\end{tabular}}
\caption{Example of a target structure being split into variables. Each color is one variable, and the positions in a variable are chosen together.}
\label{fig:structureplot}
\end{figure}


To assign initial probabilities to these variables, we follow the targeted initialization used
in SamplingDesign~\citep{tang2026}. This initialization favors A at unpaired positions and CG or GC
at paired positions, with the parameter $\epsilon\in[0,1]$ controlling the strength of this preference.
For an unpaired variable, the initial probabilities are
\begin{equation}
 p(b)=\begin{cases}
 \epsilon+(1-\epsilon)/4,&b=\mathrm{A},\\
 (1-\epsilon)/4,&b\in\{\mathrm{C,G,U}\}.
 \end{cases}
\end{equation}
\begin{table}[H]
\centering\small
\caption{Unpaired variable probabilities for $\epsilon=0.75$.}
\label{tab:unpairedinit}
\begin{tabular}{@{}lc@{}}
\toprule
Nucleotide & Probability\\
\midrule
A & 0.8125\\
C & 0.0625\\
G & 0.0625\\
U & 0.0625\\
\bottomrule
\end{tabular}
\end{table}


For a paired variable, the initial probabilities are
\begin{equation}
 p(s)=\begin{cases}
 (1+2\epsilon)/6,&s\in\{\mathrm{CG,GC}\},\\
 (1-\epsilon)/6,&\text{otherwise, }s\in\mathcal{C}.
 \end{cases}
\end{equation}
\begin{table}[H]
\centering\small
\caption{Paired variable probabilities for $\epsilon=0.75$.}
\label{tab:pairedinit}
\begin{tabular}{@{}lc@{}}
\toprule
Base pair & Probability\\
\midrule
CG & 0.4167\\
GC & 0.4167\\
AU & 0.0417\\
UA & 0.0417\\
GU & 0.0417\\
UG & 0.0417\\
\bottomrule
\end{tabular}
\end{table}

Mismatch and trimismatch variables are initialized by multiplying the unpaired probabilities of their bases.

\subsection{Sampling, scoring, and elite weighting}
\begin{figure*}[!b]
\centering
\resizebox{\textwidth}{!}{%
\begin{tikzpicture}[
  box/.style 2 args={draw=#1!75,fill=#1!10,rounded corners=3pt,line width=0.8pt,
    minimum width=#2,minimum height=1.3cm,font=\sffamily,align=center},
  flow/.style={-{Stealth[length=3.6mm]},line width=1.6pt,black!65}
]
\draw[dashed,black!45,line width=0.8pt,rounded corners=8pt,fill=black!5]
  (3.2,-1.7) rectangle (20.8,1.95);
\node[font=\sffamily\bfseries\itshape\large,text=black!65] at (12.0,1.5) {Iteration loop};
\node[box={titleblue}{2.8cm}]     (init) at (0,0)    {\large\bfseries Initialize\\[1pt]\small distribution};
\node[box={varfourcolor}{2.7cm}]  (gen)  at (4.9,0)  {\large\bfseries Sample\\[1pt]\small $N$ sequences};
\node[box={varthreecolor}{2.7cm}] (eval) at (8.4,0)  {\large\bfseries Score\\[1pt]\small $P(y\mid x)$};
\node[box={varfivecolor}{2.7cm}]  (sel)  at (11.9,0) {\large\bfseries Select elites\\[1pt]\small top $\rho$ fraction};
\node[box={varonecolor}{2.7cm}]   (wt)   at (15.4,0) {\large\bfseries Weight elites\\[1pt]\small by NED};
\node[box={vartwocolor}{2.7cm}]   (upd)  at (18.9,0) {\large\bfseries Update\\[1pt]\small $\alpha=0.5$};
\node[box={titleblue}{2.8cm}]     (ret)  at (26.0,0) {\large\bfseries Return best\\[1pt]\small sequence};
\draw[flow] (init) -- (gen);
\draw[flow] (gen) -- (eval);
\draw[flow] (eval) -- (sel);
\draw[flow] (sel) -- (wt);
\draw[flow] (wt) -- (upd);
\draw[flow] (upd.south) -- ++(0,-0.95) -| (gen.south);
\draw[flow] (upd) -- node[above=1.5mm,font=\sffamily\bfseries\large,text=black!65]{converged} (ret);
\node[font=\sffamily\bfseries\large,text=black!65] at (4.9,1.0) {+0.2\% uniform};
\end{tikzpicture}}
\caption{Overview of the EliteDesign algorithm.}
\label{fig:loop}
\end{figure*}

At each round, EliteDesign samples $N$ sequences from the current distribution with a
small amount of uniform mixing, so that every variable choice remains possible:
\begin{equation}
 p_{\mathrm{sampling}}=(1-\delta)p_{\mathrm{current}}+\delta U.
 \label{eq:sampling}
\end{equation}
In this equation, $U$ is the uniform distribution for a variable and $\delta$
is the uniform mixing rate. Each sequence in the sample is scored with
LinearPartition~\citep{zhang2020} and cached to avoid scoring the same sequence twice.
The best-scoring sequence is also tracked separately and, at the end of the
search, is rescored with ViennaRNA for accuracy~\citep{lorenz2011}.

Sequences in a sample are sorted by target probability and the top fraction
$\rho$ are selected as elites:
\begin{equation}
 N_{\mathrm{elite}}=\lceil\rho N\rceil.
\end{equation}

Each elite is then given a weight based on its NED:
\begin{equation}
 w_{\mathrm{elite}}(x)=\frac{1-\mathrm{NED}(x)}
 {\sum_{j=1}^{N_{\mathrm{elite}}}(1-\mathrm{NED}(x_j))}.
 \label{eq:weights}
\end{equation}

\subsection{Distribution updates and stopping criteria}
From the weighted elites we get the distribution $p_{\mathrm{elites}}$, which we combine with
the current distribution $p_{\mathrm{current}}$ using the smoothed cross-entropy update~\citep{deboer2005}:
\begin{equation}
 p_{\mathrm{updated}}=(1-\alpha)p_{\mathrm{current}}
 +\alpha p_{\mathrm{elites}}.
 \label{eq:update}
\end{equation}
Here $\alpha$ is the elite update rate, which is set to 0.5 by
default in our implementation. At each round we track the sequence with the
highest probability seen so far:
\begin{equation}
 P_{\mathrm{best}}=\max_{x\text{ sampled so far}}P(y\mid x).
\end{equation}
We also track the moving average over $h$ rounds:
\begin{equation}
 P_{\mathrm{avg}}=\frac{1}{h}\sum_{r=1}^{h}P_r,
\end{equation}
where $P_r$ is the average target probability in round $r$.
We stop when both of these conditions hold:
\begin{enumerate}
\item The best probability has not improved for $h$ rounds.
\item The moving average has not improved for $h$ rounds.
\end{enumerate}
\begin{table}[H]
\centering\small
\caption{Default settings of EliteDesign.}
\label{tab:defaults}
\begin{tabular}{@{}lll@{}}
\toprule
Setting & Symbol & Default\\
\midrule
Samples per round & $N$ & 500\\
Elite fraction & $\rho$ & 0.1\\
Elite update rate & $\alpha$ & 0.5\\
Uniform mixing & $\delta$ & 0.002\\
Initialization bias & $\epsilon$ & 0.75\\
Stopping condition & $h$ & 25\\
\bottomrule
\end{tabular}
\end{table}

\subsection{Algorithm summary}
Algorithm~\ref{alg:elite} summarizes EliteDesign.

\begin{algorithm}[H]
\caption{Elite-guided RNA design}
\label{alg:elite}
\small
\begin{algorithmic}[1]
\State \textbf{Input:} target $y$, initial distribution
\State $x^\star\gets\mathrm{None}$
\For{$r=1,\ldots,R_{\max}$}
    \State Sample $N$ sequences
    \State Score the sequences
    \State Update $x^\star$ if any sequence has higher $P(y\mid x)$
    \State Select $N_{\mathrm{elite}}$ sequences
    \State Weigh elites
    \State Update $p_{\mathrm{current}}$
    \If{stopping criterion is met}
        \State \textbf{break}
    \EndIf
\EndFor
\State Validate $x^\star$ with ViennaRNA
\State \Return $x^\star$
\end{algorithmic}
\end{algorithm}

\FloatBarrier
